\BeleskeID{02-s006}
% element: 02-s006:e01
% element: 02-s006:e02
% element: 02-s006:e03
% element: 02-s006:d01

Za određivanje važećeg izlaza nije dovoljno sačekati odziv samo prvog kola. Promena se prenosi kroz celu mrežu, pa se čekaju i međusignali od kojih završni izlaz zavisi. U pojednostavljenom modelu sa poznatim kašnjenjima vreme prostiranja duž jedne putanje procenjuje se sabiranjem kašnjenja kola na toj putanji.

Prelazna vrednost koja se pojavi usput ne mora biti vrednost funkcije ni za početnu ni za završnu kombinaciju ulaza. Kada različite grane prenose promenu različitom brzinom, na završno kolo kratko stiže kombinacija međusignala koja ne odgovara nijednom od ta dva stabilna stanja. Zato nije dovoljno proveriti samo stabilne vrednosti iz funkcionalne tabele; pri sintezi posebno razmatramo privremena neželjena stanja tokom prenosa promene kroz mrežu.

